% !PRINTED 662
\weberpage{662}
Es ist zu beweisen, da\ss\ dieser Grad \gesperrt{nicht kleiner als $n_1$ sein
kann}, oder was damit gleichbedeutend ist, da\ss\ die $\eta_1, \eta_2,
\ldots, \eta_{n_1}$ rational durch $\tau_1$ darstellbar sind. Dies
ergibt sich, wenn man (10) nach einem der $t_\nu$ partiell
differentiiert:
\[
\Phi'(\tau_1)\,\eta_\nu + \frac{\partial A_1}{\partial t_\nu}\tau_1^{m-1}
+ \cdots + \frac{\partial A_{m-1}}{\partial t_\nu}\tau_1
+ \frac{\partial A_m}{\partial t_\nu} = 0.
\]

Da nun $\Phi'(\tau_1)$ nicht verschwindet, wenn $m$ der m\"oglichst
niedrige Grad der Gleichung (10) ist, so erh\"alt man daraus
unmittelbar die gesuchte Darstellung von $\eta_\nu$. Also mu\ss\ $m =
n_1$ sein. Damit ist bewiesen:

2. \gesperrt{Das Funktional $\tau_1$ gen\"ugt einer irreduziblen
Gleichung $n_1$ten Grades in $z_1$.}

Setzen wir also in (9) $r = 0, 1, 2, \ldots, n_1-1$, so ergibt sich ein
System linearer Gleichungen f\"ur $\eta_1, \eta_2, \ldots, \eta_{n_1}$,
deren Determinante ein nicht verschwindendes Funktional in $z_1$ ist,
und man kann also f\"ur die Variable $t_i$ solche konstante Werte
setzen, da\ss\ die Determinante auch dann nicht verschwindet. Geht
$\tau_1$ durch diese Substitution in $\theta_1$ \"uber, so kann man aus
dem System (9) die $\eta_1, \eta_2, \ldots, \eta_{n_1}$ rational durch
$z_1$ und $\theta_1$ ausdr\"ucken.

Damit ist aber bewiesen:

3. \gesperrt{Ist $z_1$ eine beliebige, nicht konstante Funktion in
$\Omega$, und $\theta_1$ eine zweite Funktion desselben K\"orpers, die
keiner Gleichung von niedrigerem als $n_1$ten Grade in $z_1$ gen\"ugt,
so sind alle Funktionen des K\"orpers $\Omega$, also auch $z$ und
$\theta$, rational durch $z_1$ und $\theta_1$ ausdr\"uckbar.}

Zwischen irgend zwei Funktionen $\alpha, \beta$ des K\"orpers $\Omega$
besteht immer eine Gleichung mit konstanten Koeffizienten

\weq{11}{F(\alpha, \beta) = 0,}
und unter allen m\"oglichen Gleichungen dieser Form ist eine von
m\"oglichst niedrigem Grade, sowohl in bezug auf $\alpha$ als in bezug
auf $\beta$.

\begin{center}
\rule{4cm}{0.4pt}
\end{center}
